\section{符号说明}

本文建模与计算过程中使用的主要符号及其含义如表~\ref{tab:main_symbols} 所示。

\begingroup
\footnotesize
\renewcommand{\arraystretch}{1.02}
\begin{table}[H]
\centering
\caption{主要符号说明}
\label{tab:main_symbols}
\begin{tabular}{c p{6.4cm} c c}
  \toprule
  符号 & 含义 & 类型或取值范围 & 单位 \\
  \midrule

  \multicolumn{4}{l}{\textbf{模块几何参数}}\\
  $i,j$ & 硬模块编号 & $1,2,\ldots,n$ & --- \\
  $w_i^0,h_i^0$ & 模块 $i$ 的原始宽度和高度 & 正数 & 长度单位 \\
  $r_i$ & 模块 $i$ 的旋转状态 & $\{0,1\}$ & --- \\
  $w_i,h_i$ & 模块 $i$ 旋转后的宽度和高度 & 正数 & 长度单位 \\
  $(x_i,y_i)$ & 模块 $i$ 的左下角坐标 & 二维坐标 & 长度单位 \\
  $W,H$ & 全部模块最小外接轮廓的宽度和高度 & 正数 & 长度单位 \\
  $A_{\mathrm{block}}$ & 全部硬模块面积之和 & 正数 & 面积单位 \\
  $A_{\mathrm{box}}$ & 模块布局的外接轮廓面积 & 正数 & 面积单位 \\
  $\rho$ & 外接轮廓长边与短边之比 & $[1,+\infty)$ & --- \\
  $\eta$ & 外接轮廓的面积冗余率 & $[0,+\infty)$ & --- \\
  \bottomrule
\end{tabular}
\end{table}

\begin{table}[H]
\centering
\caption*{续表~\ref{tab:main_symbols}\quad 主要符号说明}
\begin{tabular}{c p{6.4cm} c c}
  \toprule
  符号 & 含义 & 类型或取值范围 & 单位 \\
  \midrule
  \multicolumn{4}{l}{\textbf{网络与固定轮廓参数}}\\
  $\mathcal{E},e$ & 网络集合及网络编号 & 有限集合 & --- \\
  $\Gamma$ & 芯片轮廓的死区率 & $[0,+\infty)$ & --- \\
  $L_{\mathrm{box}}$ & 固定正方形芯片轮廓边长 & 正数 & 长度单位 \\
  $\boldsymbol{c}_i$ & 模块 $i$ 的几何中心坐标 & 二维坐标 & 长度单位 \\
  $(X_p,Y_p)$ & 网络引脚 $p$ 的坐标 & 二维坐标 & 长度单位 \\
  $\mathcal{P}_e$ & 网络 $e$ 所包含的引脚集合 & 有限集合 & --- \\
  $d_e$ & 网络 $e$ 中的引脚数量 & 正整数 & --- \\
  $\ell_e$ & 网络 $e$ 的半周长线长 & 非负实数 & 长度单位 \\
  $\mathcal{L}$ & 全部网络的 HPWL 总和 & 非负实数 & 长度单位 \\

  \multicolumn{4}{l}{\textbf{轮廓压缩参数}}\\
  $S$ & 问题三待搜索的正方形轮廓边长 & 正整数 & 长度单位 \\
  $S_{\mathrm{lb}}$ & 正方形轮廓边长的搜索下界 & 正整数 & 长度单位 \\
  $\Gamma(S)$ & 边长为 $S$ 时对应的死区率 & $[0,+\infty)$ & --- \\

  \multicolumn{4}{l}{\textbf{非矩形模块参数}}\\
  $C_i$ & 模块 $i$ 在局部坐标系中的实际占用区域 & 有限集合 & --- \\
  $C_i^{(\theta_i)}$ & 模块 $i$ 旋转后的占用区域 & 有限集合 & --- \\
  $\theta_i$ & 非矩形模块 $i$ 的旋转角度 & $\{0^\circ,90^\circ,180^\circ,270^\circ\}$ & $^\circ$ \\
  $\widetilde C_i$ & 模块 $i$ 旋转和平移后的实际占用区域 & 有限集合 & --- \\
  $A_{\mathrm{sum}}$ & 全部非矩形模块的实际面积之和 & 正数 & 面积单位 \\
  \bottomrule
\end{tabular}
\end{table}
\endgroup
